\chapter{The Exchange Rate and the Balance of Payments (Text Ch. 25)}
\chaptermark{Exchange Rate and Balance of Payments (Text Ch. 25)}

\section{The foreign exchange market}

\subsection{Foreign currency and the forex market}

To buy goods, services or assets produced in another country we need that country's
money. Foreign bank notes, coins and bank deposits are called \emph{foreign currency}.

\begin{definition}[Foreign exchange market]
The \emph{foreign exchange market} is the market in which the currency of one country
is exchanged for the currency of another. It is not a place but a worldwide network of
banks, dealers and traders connected electronically.
\end{definition}

\subsection{Exchange rates, depreciation and appreciation}

\begin{definition}[Exchange rate]
The \emph{foreign exchange rate} is the price at which one currency exchanges for another.
In these notes the exchange rate is the number of units of foreign currency that one
Canadian dollar buys (e.g.\ U.S.\ cents per Canadian dollar).
\end{definition}

\begin{example}
If $1$ CAD $= 120$ yen, then $1$ yen $= 1/120 \approx 0.00833$ CAD. If the Canadian dollar
trades at 0.75 U.S.\ dollars, then 1 U.S.\ dollar costs $1/0.75 \approx 1.33$ CAD.
\end{example}

\begin{definition}[Depreciation and appreciation]
\emph{Currency depreciation} is a fall in the value of one currency in terms of another.
\emph{Currency appreciation} is a rise in the value of one currency in terms of another.
If the value is changed by a government that fixes the rate, a fall is called a
\emph{devaluation} and a rise a \emph{revaluation}.
\end{definition}

\begin{important}[Pitfall]
If the Canadian dollar falls from 120 yen to 100 yen, the dollar has depreciated against
the yen and the yen has appreciated against the dollar (from 0.00833 to 0.01 CAD per
yen). Always state which currency the rate is quoted in.
\end{important}

An exchange rate is a price, the price of one currency in terms of another, and like any
price it is determined by demand and supply in a market. The demand for Canadian dollars
is the supply of other currencies, and vice versa.

\section{Demand in the foreign exchange market}

People buy Canadian dollars to buy Canadian-produced goods and services (Canadian
exports) and Canadian assets. The demand for dollars is a \emph{derived demand}. The
quantity of Canadian dollars that traders plan to buy depends on
\begin{enumerate}
  \item the exchange rate ($-$: movement along the demand curve);
  \item world demand for Canadian exports ($+$);
  \item interest rates in Canada relative to other countries ($+$);
  \item the expected future exchange rate ($+$).
\end{enumerate}

\begin{important}[Law of demand for foreign exchange]
Other things remaining the same, the higher the exchange rate, the smaller is the
quantity of Canadian dollars demanded.
\end{important}

The exchange rate influences the quantity of dollars demanded for two reasons:
\begin{itemize}
  \item \textbf{Exports effect}: the larger the value of Canadian exports, the larger the
        quantity of Canadian dollars demanded. The lower the exchange rate, the cheaper
        Canadian goods are to foreigners, so the greater the value of exports and the
        quantity of dollars demanded.
  \item \textbf{Expected profit effect}: the larger the expected profit from holding
        Canadian dollars, the greater the quantity demanded today. For a given expected
        future exchange rate, the lower the exchange rate today, the larger the expected
        profit from buying dollars today and holding them.
\end{itemize}

\section{Supply in the foreign exchange market}

People sell Canadian dollars (buy foreign currency) to buy imports and foreign assets.
The quantity of Canadian dollars supplied depends on
\begin{enumerate}
  \item the exchange rate ($+$: movement along the supply curve);
  \item Canadian demand for imports ($+$);
  \item interest rates in Canada relative to other countries ($-$);
  \item the expected future exchange rate ($-$).
\end{enumerate}

\begin{important}[Law of supply of foreign exchange]
Other things remaining the same, the higher the exchange rate, the greater is the
quantity of Canadian dollars supplied.
\end{important}

\begin{itemize}
  \item \textbf{Imports effect}: the higher the exchange rate, the cheaper foreign goods
        are to Canadians, so the greater the value of Canadian imports and the larger the
        quantity of Canadian dollars supplied.
  \item \textbf{Expected profit effect}: for a given expected future exchange rate, the
        higher the exchange rate today, the larger the expected profit from selling
        Canadian dollars today (and buying them back later), so the greater the quantity
        supplied.
\end{itemize}

\section{Market equilibrium and exchange rate fluctuations}

The equilibrium exchange rate makes the quantity of Canadian dollars demanded equal the
quantity supplied. The forex market is constantly pulled to equilibrium by traders.

\subsection{Changes in the demand for Canadian dollars}

A change in any influence on the quantity of dollars people plan to buy, other than the
exchange rate, shifts the demand curve:
\begin{itemize}
  \item \textbf{World demand for Canadian exports} ($+$): if world demand for Canadian
        exports (oil, lumber, autos) increases, the demand for dollars increases.
  \item \textbf{Canadian interest rate differential} ($+$): the Canadian interest rate
        minus the foreign interest rate. A rise makes Canadian assets more attractive and
        increases the demand for dollars.
  \item \textbf{The expected future exchange rate} ($+$): if it rises, demand shifts
        rightward; if it falls, leftward.
\end{itemize}

\subsection{Changes in the supply of Canadian dollars}

\begin{itemize}
  \item \textbf{Canadian demand for imports} ($+$): rises with Canadian income, so supply
        shifts rightward.
  \item \textbf{Canadian interest rate differential} ($-$): a rise means Canadians move
        fewer funds abroad, so supply decreases.
  \item \textbf{The expected future exchange rate} ($-$): if it rises, people hold on to
        dollars, so supply decreases.
\end{itemize}

\subsection{Changes in the exchange rate}

\begin{itemize}
  \item Demand increases, supply unchanged: the exchange rate rises (appreciation).
  \item Demand decreases, supply unchanged: the exchange rate falls (depreciation).
  \item Supply increases, demand unchanged: the exchange rate falls.
  \item Supply decreases, demand unchanged: the exchange rate rises.
\end{itemize}
Because a change in the interest rate differential or in the expected future exchange rate
shifts both curves (in opposite directions), the exchange rate can change sharply even
though the quantity traded changes little. This is why exchange rates are so volatile.

\begin{center}
\begin{tikzpicture}[scale=0.8]
  \draw[->] (0,0) -- (8.5,0) node[right] {Quantity of CAD (billions/day)};
  \draw[->] (0,0) -- (0,6) node[above] {Exchange rate (US\$ per CAD)};
  \draw[thick, blue] (2,5) -- (6,1) node[below right] {$D_0$};
  \draw[thick, blue, dashed] (4,5) -- (7,2) node[below right] {$D_1$};
  \draw[thick, red] (2,1) -- (6,5) node[right] {$S$};
  \fill (4,3) circle (2.5pt);
  \fill (5,4) circle (2.5pt);
  \draw[dotted] (0,3) node[left] {0.80} -- (4,3) -- (4,0) node[below] {100};
  \draw[dotted] (0,4) node[left] {0.90} -- (5,4) -- (5,0) node[below] {125};
\end{tikzpicture}
\end{center}

\begin{example}[Solving for the exchange rate]
Let $e$ be US\$ per CAD and $Q$ billions of CAD per day, with
\[
  Q_D = 300 - 250e, \qquad Q_S = 250e - 100.
\]
Equilibrium: $300 - 250e = 250e - 100$, so $500e = 400$, $e = 0.80$ and $Q = 100$.
A world commodity boom raises demand to $Q_D = 350 - 250e$. Then $500e = 450$,
$e = 0.90$, $Q = 125$: the Canadian dollar appreciates by
$(0.90 - 0.80)/0.80 = 12.5\%$.
\end{example}

\subsection{Fundamentals, expectations and arbitrage}

\begin{definition}[Interest rate parity]
\emph{Interest rate parity} means equal rates of return on assets in different
currencies. The rate of return on a Canadian asset to a foreign holder is the Canadian
interest rate plus the expected appreciation of the Canadian dollar. In equilibrium,
\[
  \text{Canadian interest rate} - \text{foreign interest rate}
  = \text{expected depreciation of the Canadian dollar}.
\]
\end{definition}

\begin{example}
Canadian one-year rates are 2\% and U.S.\ rates are 4\%. For investors to be willing to
hold Canadian assets, the Canadian dollar must be expected to \emph{appreciate} by about
2\% over the year, so that both earn 4\%.
\end{example}

\begin{definition}[Purchasing power parity]
\emph{Purchasing power parity} (PPP) means equal value of money: a currency buys the same
basket of goods in every country. If PPP holds, the exchange rate is
\[
  e_{\text{PPP}} = \frac{\text{price of basket abroad (foreign currency)}}
  {\text{price of basket in Canada (CAD)}}.
\]
\end{definition}

\begin{example}[Big Mac PPP]
A Big Mac costs CAD 6.00 in Toronto and USD 5.00 in New York. The PPP exchange rate is
$5.00/6.00 \approx 0.83$ US\$ per CAD. If the market rate is 0.75, the Canadian dollar is
\emph{undervalued} relative to PPP, and PPP predicts it will tend to appreciate in the
long run.
\end{example}

\begin{definition}[Real exchange rate]
The \emph{real exchange rate} is the relative price of foreign-produced goods and services
to Canadian-produced goods and services. With $E$ the nominal rate (foreign currency per
CAD), $P$ the Canadian price level and $P^{*}$ the foreign price level, the number of
foreign baskets that one Canadian basket buys is
\[
  RER = \frac{E \times P}{P^{*}}.
\]
\end{definition}

\begin{example}
$E = 0.75$, a Canadian basket costs CAD 100 and the U.S.\ basket costs USD 90. Then
$RER = 0.75 \times 100/90 \approx 0.83$: one Canadian basket trades for 0.83 U.S.\ baskets.
In the short run the RER moves with $E$ (prices are sticky); in the long run it is
determined by real factors (demand and supply in goods markets), and the nominal rate
adjusts to the price levels.
\end{example}

\section{Exchange rate policy}

\begin{itemize}
  \item \textbf{Flexible exchange rate}: determined by demand and supply with no direct
        intervention by the central bank. Canada has had a flexible rate since 1970 (it
        was fixed 1962--1970).
  \item \textbf{Fixed exchange rate}: pegged at a value decided by the government or
        central bank. To keep it fixed, the central bank buys its own currency (using
        foreign reserves) when there is excess supply and sells it when there is excess
        demand. It cannot defend an overvalued rate forever: reserves run out.
  \item \textbf{Crawling peg}: a rate that follows a path chosen by the central bank
        (China's yuan after 2005 is the textbook example).
\end{itemize}

\section{Financing international trade: the balance of payments}

\begin{definition}[Balance of payments accounts]
A country's \emph{balance of payments accounts} record its international trading,
borrowing and lending in three accounts.
\begin{enumerate}
  \item \textbf{Current account}: receipts from exports of goods and services, minus
        payments for imports, plus net interest income and net transfers received from
        abroad. Its balance is roughly net exports.
  \item \textbf{Capital and financial account}: foreign investment in Canada minus
        Canadian investment abroad.
  \item \textbf{Official settlements account}: the change in Canada's official
        international reserves (the government's holdings of foreign currency), recorded
        with a negative sign when reserves increase.
\end{enumerate}
The sum of the three balances is always zero:
\[
  \text{current account} + \text{capital and financial account} +
  \text{official settlements account} = 0.
\]
\end{definition}

\begin{example}
In a year (billions of dollars): exports 600, imports 620, net interest income $-20$,
net transfers $-2$. Foreign investment in Canada 150, Canadian investment abroad 110.
\begin{align*}
  \text{current account} &= 600 - 620 - 20 - 2 = -42,\\
  \text{capital and financial account} &= 150 - 110 = 40,\\
  \text{official settlements account} &= -(-42 + 40) = +2,
\end{align*}
meaning official reserves \emph{decreased} by 2 billion.
\end{example}

\begin{definition}[Borrowers, lenders, debtors, creditors]
A \emph{net borrower} borrows more from the rest of the world than it lends in a given
year (a current account deficit); a \emph{net lender} lends more than it borrows. A
\emph{debtor nation} has, over its history, borrowed more than it has lent (a stock); a
\emph{creditor nation} has invested more abroad than foreigners have invested in it.
\end{definition}

Being a net borrower is not necessarily bad: it matters whether the borrowing finances
investment (which generates income to repay it) or consumption. Canada was a debtor
nation for most of its history, borrowing to build its capital stock.

\subsection{The current account balance and national saving}

From $Y = C + I + G + NX$ and $Y = C + S + T$,
\[
  NX = (T - G) + (S - I).
\]
Net exports equal the government sector balance plus the private sector balance. A
government deficit, other things equal, tends to reduce net exports (the ``twin
deficits'' idea), but in practice the private balance moves too.

\begin{example}
The government runs a deficit of 30 ($T - G = -30$) and private saving falls short of
investment by 10 ($S - I = -10$). Then $NX = -30 - 10 = -40$: the economy runs a current
account deficit of about 40 and borrows that amount from the rest of the world.
\end{example}

\section{Summary}

\begin{itemize}
  \item The exchange rate is the price of one currency in terms of another; it is
        determined by demand and supply in the forex market.
  \item Demand for CAD: falls with the exchange rate (exports and expected profit
        effects); rises with world demand for exports, the Canadian interest rate
        differential, and the expected future exchange rate.
  \item Supply of CAD: rises with the exchange rate (imports and expected profit effects)
        and with Canadian demand for imports; falls with the interest rate differential and
        the expected future exchange rate.
  \item Interest rate parity and purchasing power parity anchor exchange rates;
        $RER = EP/P^{*}$.
  \item Exchange rate regimes: flexible (Canada), fixed, crawling peg.
  \item Balance of payments: current account $+$ capital and financial account $+$
        official settlements $= 0$; $NX = (T - G) + (S - I)$.
\end{itemize}
